\chapter*{B3\quad 试卷 3 参考答案}
\addcontentsline{toc}{chapter}{B3\quad 试卷 3 参考答案}
\markboth{B3\quad 试卷 3 参考答案}{}

\begin{tishi}
本答案\textbf{照录原文档给出的标准答案串}（打印在原文档末页），\textbf{一个字母都没有改动}；
原文档判断题用的“对勾/叉”符号，本册按规范写作\textbf{（对）}、\textbf{（错）}。
每题另附一句辨析，供选择题把四个（或五个）选项都弄懂用。
\textbf{答案册不标星级。}
\end{tishi}

\section*{一、单项选择题（30 题 $\times$ 2 分）}

\noindent\textbf{第 1--5 题}\quad 1.D\quad 2.C\quad 3.A\quad 4.D\quad 5.A\par
\noindent\textbf{第 6--10 题}\quad 6.B\quad 7.A\quad 8.B\quad 9.C\quad 10.A\par
\noindent\textbf{第 11--15 题}\quad 11.D\quad 12.D\quad 13.C\quad 14.A\quad 15.C\par
\noindent\textbf{第 16--20 题}\quad 16.D\quad 17.B\quad 18.C\quad 19.B\quad 20.A\par
\noindent\textbf{第 21--25 题}\quad 21.B\quad 22.A\quad 23.C\quad 24.B\quad 25.C\par
\noindent\textbf{第 26--30 题}\quad 26.C\quad 27.D\quad 28.B\quad 29.D\quad 30.A

\section*{二、多项选择题（10 题 $\times$ 2 分）}

\noindent 1.AB\quad 2.ABCE\quad 3.BCDE\quad 4.ABCD\quad 5.ABD\par
\noindent 6.ABCD\quad 7.AB\quad 8.AB\quad 9.CDE\quad 10.AB

\section*{三、判断题（20 题 $\times$ 1 分）}

\noindent\textbf{第 1--5 题}\quad 1.（对）\quad 2.（错）\quad 3.（对）\quad 4.（错）\quad 5.（错）\par
\noindent\textbf{第 6--10 题}\quad 6.（错）\quad 7.（对）\quad 8.（对）\quad 9.（对）\quad 10.（错）\par
\noindent\textbf{第 11--15 题}\quad 11.（错）\quad 12.（对）\quad 13.（对）\quad 14.（对）\quad 15.（错）\par
\noindent\textbf{第 16--20 题}\quad 16.（错）\quad 17.（对）\quad 18.（错）\quad 19.（错）\quad 20.（错）

\section*{四、逐题辨析（选择、判断）}

\noindent\textbf{单选第 1 题（流体的定义）}\quad
“流体”的定义落点在\textbf{受剪切力时不能保持静止}（选项 D），
这正是流体与固体的分界。选项 A 说的是气体、B 说的是理想化的不可压缩流体、
C 说的是理想（无黏）流体，三个都是“某一类流体”的特征，不是流体的定义。
注意：原文档题干把它写成了“牛顿流体是……”，与四个选项矛盾，应按“流体是……”读。

\noindent\textbf{单选第 2 题、单选第 9 题（连续性方程）}\quad
连续性方程是\textbf{质量守恒}在流体上的数学表达（2 题选 C）。
不可压缩流体恒定总流的连续性方程即“流量沿程不变”：
\[ v_1A_1=v_2A_2 \]
所以 9 题选 C。若流体可压缩，则须写成 $\rho_1v_1A_1=\rho_2v_2A_2$，选项里没有这一项，
这也是第 9 题按“不可压缩”理解的原因。

\noindent\textbf{单选第 3 题、第 15 题（沿程阻力系数 $\lambda$）}\quad
圆管层流 $\lambda=64/\Rey$，\textbf{只与雷诺数有关}（3 题选 A）；
紊流分三个区：水力光滑区 $\lambda$ 只与 $\Rey$ 有关、过渡区 $\lambda$ 与 $\Rey$ 和相对粗糙度都有关、
粗糙区（阻力平方区）$\lambda$ 只与相对粗糙度有关。15 题问“紊流过渡区”，故选 C。

\noindent\textbf{单选第 4 题（动量方程）}\quad
动量方程是\textbf{矢量方程}，流速和作用力都是有方向的量，列式时必须写成分量式（选 D）。
2015 年计算第 4 题（水平弯管求 $F_x$、$F_y$）与 2026 年回忆版计算第 3 题（教材习题 4.12）
考的都是这一点。

\noindent\textbf{单选第 5、6 题与多选第 2 题（雷诺数）}\quad
$\Rey=\dfrac{vd}{\nu}$，是\textbf{惯性力与黏滞力之比}（5 题选 A）。
由表达式可见 $\Rey$ 与\textbf{平均流速 $v$、运动黏滞系数 $\nu$、管径 $d$} 有关，
与流体的比重、静压力无关，故多选 2 题选 ABCE。
圆管水流的下临界雷诺数为 \textbf{2000}（6 题选 B）——上临界雷诺数约为 4000，
两条临界值不要混。

\noindent\textbf{单选第 8 题（黏度换算）}\quad
直接用 $\mu=\rho\nu$：
\[ \mu=800\times3\times10^{-6}=2.4\times10^{-3}\ \mathrm{Pa\cdot s} \]
选 B。注意 $800\ \mathrm{kg/m^3}$ 对应的是油类，不是水。

\noindent\textbf{单选第 10 题（水力半径）}\quad
水力半径定义为过流断面面积与湿周之比 $R=A/\chi$。矩形断面渠道
$A=bh=2\times1=2\ \mathrm{m^2}$，湿周 $\chi=b+2h=2+2\times1=4\ \mathrm{m}$，
故 $R=2/4=0.5\ \mathrm{m}$，选 A。

\noindent\textbf{单选第 7、13 题与判断第 9 题（长管的串联与并联）}\quad
\textbf{串联}管道无支流分出，各管段\textbf{流量相等}（7 题选 A），水头损失沿程叠加
$h_f=\sum h_{fi}$；\textbf{并联}管道各支管的两端是同一对节点，
故\textbf{沿程水头损失相等}（13 题选 C），流量则按各支管阻抗分配。
长管计算中把局部损失和出口动能忽略不计，认为水头\textbf{全部消耗在沿程水头损失上}，
判断第 9 题判“对”。

\noindent\textbf{单选第 11 题（流线与迹线）}\quad
\textbf{定常流动}时流线与迹线重合（选 D）。无旋、均匀、平面都是对流动的另外的分类，
与“流线是否与迹线重合”无关。

\noindent\textbf{单选第 12 题（水击）}\quad
减轻水击的常用措施是：\textbf{加大管径}（降低流速）、\textbf{缩短管道长度}、
\textbf{延长阀门的启闭时间}、设置调压塔或空气阀等。
“\textbf{减少阀门的启闭时间}”恰好相反，会加剧水击，故它是“不能采取的措施”，选 D。
（原文档截图此处另有读者手写的“延长”二字，属批注，不是原题文字。）

\noindent\textbf{单选第 14 题（测压管水头线）}\quad
仅受重力作用的静止液体，各点测压管水头 $z+p/\gamma$ 相同，
测压管水头线是\textbf{水平线}（选 A）。这与 2022 年单选第 3 题
（测压管水头线坡度小于零时沿程下降）是一对正反设问，务必对照记住。

\noindent\textbf{单选第 16 题（绝对压强）}\quad
绝对压强等于大气压加水柱压强：
\[ p_{abs}=p_a+\rho gh=98+9.8\times10=196\ \mathrm{kPa} \]
选 D（$196\ \mathrm{kPa}$ 正好是两个大气压）。选项 C 的 $147\ \mathrm{kPa}$ 相当于 15 m 水柱，
数值与 10 m 不符，是干扰项。注：$9.8\times10=98\ \mathrm{kPa}$，正好是一个大气压。

\noindent\textbf{单选第 17 题、判断第 7 题与多选第 7 题（压强基准）}\quad
绝对压强恒为正（选 17 题 B、D 的前半句对、后半句错，故 17 题选 B）；
相对压强以当地大气压为零点，\textbf{可正、可负、可为零}，所以判断第 7 题判“对”。
三者关系为
\[ p_{abs}=p+p_a,\qquad p=p_{abs}-p_a,\qquad p_v=p_a-p_{abs} \]
多选第 7 题原卷给 AB（即 $p_{abs}=p_v+p_a$、$p=p_{abs}-p_a$），
说明原卷把 $p_v$ 当“与相对压强同量值”的记号使用；若严格按“真空度”理解，则 C 也成立。
考场上以本题题干给出的记号习惯为准。

\noindent\textbf{单选第 18 题（均匀流场）}\quad
\textbf{均匀流}指流场中各空间点的运动要素不随空间位置变化；
题中“每一空间点上物理参数均相等”说的正是空间上均匀，故选 C。
与之对照，\textbf{定常流}说的是同一空间点的参数不随时间变化；两者一个管“空间”、一个管“时间”。

\noindent\textbf{单选第 19 题（动能修正系数 $\alpha$）}\quad
$\alpha$ 是断面上实际动能与按平均流速计算的动能之比。
圆管层流流速呈抛物线分布，$\alpha=2$（选 B）；紊流流速分布较均匀，$\alpha\approx1$。
2015 年填空“层流断面平均速度等于 0.5 倍最大速度”与本题是同一结论的两个侧面。

\noindent\textbf{单选第 20 题（层流沿程损失）}\quad
层流时 $\lambda=64/\Rey\propto1/v$，代入 $h_f=\lambda\dfrac{l}{d}\dfrac{v^2}{2g}$ 得
$h_f\propto v^{1}$，即\textbf{1 次方}（选 A）。
紊流水力光滑区 $h_f\propto v^{1.75}$、阻力平方区 $h_f\propto v^{2}$——
选项里的 1.75、1.75$\sim$2、2 都是这两个区的指数，正是干扰项，
所以“层流＝1 次方”必须记牢。

\noindent\textbf{多选第 1 题（描述流体运动的方法）}\quad
只有\textbf{欧拉法}与\textbf{拉格朗日法}两种（选 AB）。“分析法”不是描述流体运动的方法。

\noindent\textbf{多选第 3 题（运动参数）}\quad
描写流体运动状态的运动参数主要是\textbf{速度、加速度、角速度}，原卷答案取 BCDE
（把重力加速度也计入）；平均速度是统计量，不是运动参数。按原卷答案记 BCDE。

\noindent\textbf{多选第 8 题（管道的分类）}\quad
按水头损失计算方法的不同，管道分为\textbf{长管}与\textbf{短管}（选 AB）。
圆管、粗糙管、光滑管是按几何形状或壁面状况分类，不是按计算方法分类。

\noindent\textbf{多选第 9 题（流体的物理特性）}\quad
除密度、比容外，流体的主要物理特性还有\textbf{黏性、压缩性、重度}（选 CDE）。
体积、重量是几何量与力学量，不是流体的物理特性。

\noindent\textbf{判断第 11 题（连续性方程的来源）}\quad
总流连续性方程是\textbf{质量守恒}（质量守恒定律）在总流上的表达，
不是从能量守恒方程推导出来的，故判“错”。能量守恒导出的是伯努利方程。

\noindent\textbf{判断第 12、13 题（边界层与水力光滑管）}\quad
紊流边界层的定义就是边界层内流体以\textbf{紊流运动为主}，判“对”（第 12 题）；
\textbf{层流底层厚度大于管壁粗糙度}时，粗糙峰完全被层流底层淹没、对流动没有影响，
这种管叫\textbf{水力光滑管}，判“对”（第 13 题）。反过来则是水力粗糙管。

\noindent\textbf{判断第 14、19、20 题（黏性相关）}\quad
$\nu=\mu/\rho$：只说“动力黏度相等”而密度不同，运动黏度就不相等，第 14 题判“对”的前提是
\textbf{同一种不可压缩液体}（密度相同）；常温常压下\textbf{空气}的运动黏度
（约 $1.5\times10^{-5}\ \mathrm{m^2/s}$）比水（约 $1.0\times10^{-6}\ \mathrm{m^2/s}$）\textbf{大}，
故第 19 题判“错”；由牛顿内摩擦定律 $\tau=\mu\dfrac{\dd u}{\dd y}$，
切应力不只取决于流速梯度，还取决于该液体的动力黏度，故第 20 题判“错”。

\noindent\textbf{判断第 15、16、17 题（压缩性与黏性）}\quad
体积模量 $K$ 越大表示越\textbf{难}压缩，第 15 题判“错”；
液体的黏性\textbf{只在运动（有相对运动）时表现为切应力}，静止液体不承受切应力，
但“静止时也具有黏滞性”这一性质判断，第 16 题判“错”、第 17 题判“对”。

\noindent\textbf{判断第 18 题（尼古拉兹试验）}\quad
原卷判“错”。尼古拉兹试验用人工均匀砂粒粗糙管，直接测定的对象是
\textbf{沿程阻力系数 $\lambda$} 随雷诺数与相对粗糙度 $\Delta/d$ 的变化规律；
题干把它写成“沿程水头损失”，用词不符，故判错。
若题目写成“研究沿程阻力系数 $\lambda$……”，则应判对，注意题干用词。

\section*{五、超纲题（流体机械）提示}

本卷有 22 题属\textbf{泵与风机（流体机械）}，对应赵琴教材第 11--14 章，
\textbf{不在 818 初试范围}（初试只考前 8 章），答案照录如下，供复试参考：
\begin{enumerate}
  \item 单选第 21--30 题：21.B（平衡盘平衡轴向力）、22.A（平衡盘能自动调整轴向力）、
        23.C（口环是叶轮与泵壳间的密封）、24.B（轴流式流量大）、
        25.C（工况点落在最高效率点附近最经济）、26.C（叶轮使流体获得能量）、
        27.D（离心式）、28.B（利用升力）、29.D（收集流体并把部分动能变为压能）、
        30.A（螺旋线）。
  \item 多选第 4、5、6、10 题：4.ABCD（水力损失、容积损失、机械损失），
        5.ABD（吸水段、排水段、中段），6.ABCD（200D43$\times$3 型号含义），
        10.AB（降低声源噪声、传输路径噪声控制）。
  \item 判断第 1--6、8、10 题：1.对、2.错、3.对、4.错、5.错、6.错、8.对、10.错。
\end{enumerate}
